%!TEX root = ../template.tex
%%%%%%%%%%%%%%%%%%%%%%%%%%%%%%%%%%%%%%%%%%%%%%%%%%%%%%%%%%%%%%%%%%%%
%% chapter6.tex
%% NOVA thesis document file
%%
%% Chapter with lots of dummy text
%%%%%%%%%%%%%%%%%%%%%%%%%%%%%%%%%%%%%%%%%%%%%%%%%%%%%%%%%%%%%%%%%%%%

\typeout{NT FILE chapter6.tex}%

\chapter{OCAML-FLAT - Implementation}
\label{cha:implementation}

OCaml-FLAT is one of the two main components that are part of the existing structure of this project. It consists of the core logic of the tool and 
is implemented entirely in OCaml. In this chapter, we will address OCaml-FLAT, which contains the main operations implemented for attribute grammars.

OCaml-FLAT is a library composed of several modules, and each module is built as a class. This library, up to the moment of writing this document, implements 
finite automata, regular expressions, context-free grammars, pushdown automata, Turing machines and unrestricted grammars. as well as some supporting modules, 
such as operations on JSON syntax, error management, parsers, among others.

For the implementation of attribute grammars in OCaml-FLAT, the central goal is to calculate synthesized and inherited attributes, validate grammar, and perform cycle detection. The library represents an attribute grammar as a set of rules with equations and conditions attached to each production, and it 
enforces a good construction before any evaluation can proceed. Validation checks range from context-free sanity (reusing and extending CFG validation) to attribute 
typing, direction constraints (inherited vs.\ synthesized), and global acyclicity of attribute dependencies. Only after this pipeline succeeds do we 
execute the attribute equations over a parse tree.

Two aspects are key to the design. First, expression typing and equation validation ensure each assignment is meaningful and referentially correct. 
Second, the evaluation order guarantees that inherited attributes of a child are computed before descending into that child, and synthesized attributes 
of the head are computed after its children return---yielding a left-to-right. Together, these choices produce a robust and explainable evaluator.

Typing proceeds by a straightforward rule: attributes are associated to a base type inferred from their name---by default integers, with \texttt{"s"} 
denoting string and \texttt{"b"} denoting boolean---while constants and operators are type checked accordingly. Binary operators are admitted only when 
operand types align (e.g., \texttt{+} on \texttt{int} or \texttt{string}, \texttt{-}, \texttt{*}, \texttt{/} on \texttt{int}, comparisons across equal 
types), and any mismatch is flagged at validation time. Applications of attributes use \texttt{(attr, (var, i))}, for humans, \texttt{v(S0)} , where \texttt{i} identifies the head 
or a specific child occurrence.

On the left-hand side of equations, the direction discipline is enforced: assignments targeting the head must write synthesized attributes, whereas 
assignments targeting a child must write inherited attributes. Indices are checked against the rule's body to prevent out-of-range references, and 
all variables must belong to the grammar's variable set. If types on both sides of an equation do not match, validation fails. Conditions are validated 
separately and must type check to boolean.

At evaluation time, the library models a local environment as a list of nodes, headed by the current symbol and followed by its children. Each node 
carries the set of evaluated attributes for that symbol occurrence. Attribute applications resolve by locating the appropriate node---head or child 
occurrence---and reading the requested attribute value. Updates are performed by replacing the targeted attribute binding on the designated node, 
preserving immutability of the structure. A stable ordering of equations at the head and per child ensures determinism of side-effect-free updates 
within a node.

The evaluation strategy is left-to-right on children. For each child, all inherited equations targeting that child are applied in the head's environment 
before the recursive descent. The child is then evaluated. Once all children have been processed, equations that target the head (i.e., synthesized 
attributes) are evaluated, and node-local conditions are checked. This yields a two-phase per-node schedule---child inherited first, head synthesized 
later---that matches the intuition for attribute flow.

Cycle detection is guarded by a dependency graph construction. For every equation, the implementation identifies the destination attribute occurrence 
(head or specific child) and the set of attribute references appearing on the right-hand side. From these, it builds a directed graph whose nodes 
are ``(symbol, occurrence, attribute)'' triplets and whose edges represent ``RHS depends on LHS.'' A depth-first search detects cycles; if any cycle 
exists, the grammar is rejected as having circular attribute dependencies. This check is integrated into the overall \texttt{validate} routine, 
preventing non-terminating evaluations at runtime.

The attribute grammar layer interacts with the context-free layer. Conversions are provided in both directions: an attribute grammar can be projected to 
a CFG by erasing attributes, and a CFG can be lifted to an attribute grammar with empty attribute sets and equation blocks. Word acceptance may, when 
appropriate, delegate to the CFG acceptance procedure, while a separate routine attempts attribute evaluation on a given parse tree, reporting success 
or failure. A generator uses the CFG projection to enumerate terminal strings up to depth and count limits, which is useful for testing attribute 
specifications independently of a parser.

Reference semantics for attribute occurrences are explicit. Index \texttt{0} denotes the head; index \texttt{-1} is normalized to the head when the variable 
equals the head, and otherwise to the first occurrence of the named child; positive indices select the \(k\)-th matching child. Occurrence computations and 
equation ordering are derived from the rule body, ensuring that references are stable and reproducible even when the same child symbol appears multiple 
times. This removes ambiguity at both validation and evaluation time.

\section{Data types created for supporting attribute grammars} 

The main data types defined to simplify the manipulation of attribute grammars are: \texttt{t}, which represents an attribute grammar (alphabet, 
variables, synthesized and inherited attributes, initial variable, and set of rules); \texttt{rule}, which models a production with \emph{head}, \
emph{body}, a set of \texttt{equations}, and a set of \texttt{conditions}; and \texttt{equation}, which represents an assignment between 
\texttt{expression}s. Additionally, \texttt{expression} and \texttt{value} describe the expressions and constants used in equations; 
\texttt{condition} represents boolean predicates; and \texttt{attribute}, \texttt{attributes}, and \texttt{attrArg} handle the identification of 
attributes and their occurrences.

\begin{verbatim}
type attribute = symbol
    type attributes = attribute set
    type attrArg = variable * int
    type value =
        | Int of int
        | String of string
        | Bool of bool
    type expression =
        | Const of value
        | Apply of attribute * attrArg (* l(A2) *)
        | Expr of string * expression * expression
    type equation = expression * expression
    type equations = equation set
    type condition = expression 
    type conditions = condition set

    type rule = {
        head : variable;
        body : word;
        equations : equations;
        conditions : conditions
    }
    type rules = rule set

    type t = {
        alphabet : symbols;
        variables : variables;
        synthesized : attributes;
        inherited : attributes;
        initial : variable;
        rules : rules
    }
\end{verbatim}


\subsection*{Notes on \texttt{attrArg} and \texttt{Apply}}

\texttt{attrArg} = \texttt{(variable * int)} Identifies the target node of the current rule:
\begin{itemize}
  \item \textbf{0} = head of the rule; \quad \textbf{$k>0$} = $k$-th occurrence in the body; \quad \textbf{$-1$} = default (head if the variable is the head, otherwise first occurrence).
\end{itemize}

\texttt{Apply} = \texttt{(attribute * attrArg)} Reads/writes attribute \texttt{attr} on that node:
\begin{itemize}
  \item RHS = read; \quad LHS = write (\textbf{head} $\Rightarrow$ synthesized, \textbf{child} $\Rightarrow$ inherited).
\end{itemize}

Example
\texttt{E -> E * F \{ v(E0) = v(E1) * v(F) \}} uses \texttt{E0=(E,0)}, \texttt{E1=(E,1)}, \texttt{F=(F,-1)$\Rightarrow$(F,1)}.




\section{The \texttt{calcAttributes} function} 

\texttt{calcAttributes} evaluates all attributes in a parse tree according to a given attribute grammar. 
It ensures that each child node receives its inherited attributes before that child is evaluated, and that the parent node computes its synthesized attributes 
only after all children have been evaluated. 
The function returns a tree with the same shape, but with attribute maps updated on the head node and on every descendant. 


\begin{verbatim}
let rec calcAttributes (ag : t) (pt : parseTree) : parseTree =
    match pt with
    | Leaf _ -> pt
    | Node (((head_sym, _) as head), children) ->
        let rule = getRootRule ag (Node (head, children)) in
        let children' = eval_children_left_to_right ag head rule.equations children in
        let head_eqs = set_filter (eq_targets_head head_sym) rule.equations in
        let env2 = head :: List.map getRoot children' in
        let env3 = apply_equations_at_head head_sym env2 head_eqs in
        check_conditions_at_node head_sym rule.conditions env3;

        match env3 with
        | (new_head_sym, new_head_evs) :: child_envs ->
            let children'' = zip_update_children children' child_envs in
            Node ((new_head_sym, new_head_evs), children'')
        | [] ->
            failwith "calcAttributes: empty environment after synthesized equations"
\end{verbatim}


\subsection{Inputs and Output}
The function takes an attribute grammar \texttt{ag} and a parse tree \texttt{pt}. If \texttt{pt} is a leaf, it is returned unchanged. 
If \texttt{pt} is an internal node, the function computes attributes for the node and for all its children, returning a structurally identical tree
whose nodes carry the evaluated attribute bindings. 

\subsection{Evaluation Procedure at an Internal Node}
Let the current node be \texttt{Node((head\_sym, head\_evs), children)}.
\begin{enumerate}
  \item \textbf{Rule identification.} The function determines the unique grammar rule whose head symbol equals \texttt{head\_sym} and whose body matches the 
sequence of child symbols at this node. If no rule matches, a diagnostic is raised that includes the encountered head/body and the candidate rules 
that share the same head. This step fixes the set of equations and conditions that apply at the node. 

  \item \textbf{Inherited attributes for children (left-to-right).} Children are processed strictly from left to right. For the current child occurrence, 
the function builds an \emph{environment} that consists of the head node, the already processed children (in order, with any attributes computed so far), 
and the remaining children (not yet processed). From the node's rule, it selects only those equations that assign \emph{inherited} attributes to this specific 
child occurrence (identified by symbol and 1-based occurrence index), and applies them to the environment in a stable order. The child's root is updated with these inherited 
values, and then \texttt{calcAttributes} is called recursively on that child to evaluate its subtree. This guarantees that a child's inherited attributes may depend on the parent 
and on earlier siblings, but not on later siblings \cite{AhoSethiUllman1986}. 

  \item \textbf{Synthesized attributes for the head.} Once all children are fully evaluated, the function gathers only the equations that target the head 
(i.e., assign \emph{synthesized} attributes). It constructs a head-first environment containing the head node followed by the now-evaluated child roots and 
applies those equations in a deterministic, stable order, updating the head's attribute bindings. 

  \item \textbf{Condition checking.} The boolean conditions attached to the rule are evaluated in the same head-first environment. Each condition must evaluate 
to \texttt{true}. If a condition fails or evaluates to a non-boolean value, the evaluation is considered unsuccessful at this node. 

  \item \textbf{Reconstruction.} Finally, the function rebuilds and returns the node by replacing the head and each child root with their updated attribute 
environments. The tree's topology is unchanged; only attribute maps are updated. 
\end{enumerate}

\subsection{Deterministic Order and Rationale}
The strict left-to-right order establishes a predictable flow of information: inherited attributes for child \( i \) may depend on the head and on children
\( 1..i-1 \), but never on \( i+1..n \). This sequencing avoids evaluation-time circularities and ensures each child observes a stable context when it is 
evaluated. A separate dependency-graph check elsewhere detects global cycles among equations. 

\subsection{Error Handling}
The surrounding code first calls \texttt{validate}; if it returns \texttt{false}, evaluation is skipped and the check reports \texttt{false}. 
If it returns \texttt{true}, structural and type errors (invalid references, out-of-range occurrences in equations, operator/type mismatches, cycles) are excluded; 
any remaining failures are runtime (e.g., division by zero), conditions that evaluate to \texttt{false}, or inconsistencies in hand-built parse trees. 
Diagnostics are reported at the point of detection (candidate rules on rule-selection mismatches; node identification on condition failures).

\section{Auxiliary Methods Used by \texttt{calcAttributes}}
This section focuses on the local worker that drives left-to-right child evaluation for \texttt{calcAttributes}. In the implementation, \texttt{calcAttributes} 
delegates the child pass to \texttt{eval\_children\_left\_to\_right}, which defines a local recursive function \texttt{loop}. 
This \texttt{loop} is the auxiliary routine that realizes the stepwise propagation of inherited attributes to each child and then triggers 
recursive evaluation of that child. 

\subsection{\texttt{eval\_children\_left\_to\_right}}

Evaluate children in strict left-to-right order, applying the inherited-attribute equations that target each child occurrence before recursively evaluating that 
child's subtree. This ensures that each child receives its required inherited attributes possibly dependent on the head and on previously evaluated siblings, 
but not on later siblings.

This function receives a grammar \texttt{ag}, the current node's head \texttt{head}, the full set of equations for the node's rule \texttt{eqs}, and the list 
of child parse trees \texttt{children}. And returns a list of fully evaluated child parse trees, in original order. 

The function computes the $(\text{symbol}, \text{occurrence})$ pair for each child (e.g., second \texttt{Expr} child). It defines a 
local recursive worker \texttt{loop} that maintains three lists: already processed children (reversed), 
remaining children, and the corresponding $(\text{symbol}, \text{occurrence})$ pairs for the remaining ones. For the current child, an evaluation 
environment is built as \texttt{head :: processed\_roots @ remaining\_roots}. Inherited equations targeting exactly this child occurrence are filtered and 
applied to the environment, producing the child's inherited attributes. The child root is updated with those attributes, the child is recursively evaluated 
via \texttt{calcAttributes}, and the loop advances to the next child. Any mismatch between children and occurrence annotations raises a clear failure.


\begin{verbatim}
and eval_children_left_to_right
    (ag : t)
    (head : node)
    (eqs : equation Set.t)
    (children : parseTree list)
    : parseTree list =
  let head_sym = fst head in
  let occs = child_occurrences children in
  let rec loop
      (processed : parseTree list)        (* evaluated children, reversed *)
      (remaining : parseTree list)        (* children yet to process     *)
      (remaining_o : (symbol * int) list) (* their (sym,occ)             *)
      : parseTree list =
    match remaining, remaining_o with
    | [], [] ->
        List.rev processed
    | child :: rest, occ :: occs_rest ->
        let processed_nodes = List.rev processed |> List.map getRoot in
        let env_before =
          head :: (processed_nodes @ (List.map getRoot (child :: rest)))
        in
        let inh_eqs = set_filter (eq_targets_child head_sym occ) eqs in
        let env_after_inh =
          apply_equations_at_head head_sym env_before inh_eqs
        in
        let child_node_after_inh =
          match env_after_inh with
          | _head_after :: env_tail ->
              pick_child_node_by_occ occ env_tail
          | [] ->
              failwith
                "calcAttributes: empty environment after inherited equations"
        in
        let child_with_inh = updateRoot child child_node_after_inh in
        let child_evaluated = calcAttributes ag child_with_inh in
        loop (child_evaluated :: processed) rest occs_rest
    | _ ->
        failwith
          "calcAttributes: internal error (children/occurrences mismatch)"
  in
  loop [] children occs
\end{verbatim}

The call to \texttt{eval\_children\_left\_to\_right} is the point where the local auxiliary function \texttt{loop} is used. 

The \texttt{loop} function above is the auxiliary worker for \texttt{calcAttributes}. 

\subsection{Role of the Auxiliary \texttt{loop} Function}
The \texttt{loop} function performs a left-to-right sweep over the node's children. For each child in order, it computes and writes that child's inherited 
attributes, recursively evaluates the child subtree, and carries the result forward before proceeding to the next child. 
This realizes the required dependency discipline: a child may depend on the head and on previously evaluated siblings, but not on later siblings. 

The function carries three accumulators:
\begin{itemize}
  \item \texttt{processed}: the list of children already evaluated; kept in reverse order for efficient cons operations. These are fully evaluated subtrees. 
  \item \texttt{remaining}: the list of children yet to process, with the next child at the head. 
  \item \texttt{remaining\_o}: the list of $(\textit{symbol}, \textit{occurrence})$ pairs corresponding one-to-one with \texttt{remaining}, used 
  to target inherited-equation assignments to the correct child occurrence. 
\end{itemize}

\textbf{Key invariants}
\begin{enumerate}
  \item \texttt{processed @ remaining} is a permutation of the original \texttt{children}, preserving relative order between the still-unseen suffix and the 
  already-evaluated prefix. The invariant is maintained by consing the newly evaluated child onto \texttt{processed} and dropping it from \texttt{remaining}. 
  \item \texttt{remaining\_o} matches \texttt{remaining} position-by-position; each pair $(s,k)$ gives the symbol and $k$-th occurrence index of the corresponding 
  child among all children of this node. This allows selection of inherited equations and of the correct child node in the environment. 
  \item Every tree in \texttt{processed} is fully evaluated (its inherited attributes were set before the recursive call, and its synthesized attributes 
  and conditions were handled inside that recursive \texttt{calcAttributes}). 
\end{enumerate}

\subsection{Step-by-Step Operation of \texttt{loop}}
For the case \verb|child :: rest, occ :: occs_rest|:
\begin{enumerate}
  \item Compute \texttt{processed\_nodes} as the roots of already-evaluated children (in original order). Then build the evaluation environment for this step:
  \[
    \texttt{env\_before} =
    \texttt{head} \;::\;
    \texttt{processed\_nodes} \; @ \;
    \texttt{(roots of (child :: rest))}.
  \]
  This environment exposes the head, earlier siblings, the current child, and later siblings to the equations to be applied. 

  \item Select only inherited-attribute equations that target the current occurrence \texttt{occ}:
  \[
    \texttt{inh\_eqs} =
    \{\, e \in \texttt{eqs} \mid \texttt{eq\_targets\_child head\_sym occ}(e)\,\}.
  \]
  Apply them in a stable order at the head of \texttt{env\_before}:
  \[
    \texttt{env\_after\_inh} =
    \texttt{apply\_equations\_at\_head head\_sym env\_before inh\_eqs}.
  \]
  The effect is to update exactly the target child’s inherited attributes in the environment. 

  \item Extract the updated root node for the current child occurrence from the environment tail using \texttt{pick\_child\_node\_by\_occ occ}, and write 
  it back into the child tree with \texttt{updateRoot}:
  \[
    \texttt{child\_with\_inh} = \texttt{updateRoot child child\_node\_after\_inh}.
  \]
  If the environment is malformed or the occurrence is not found, a descriptive failure is raised. 

  \item Recursively evaluate the child subtree:
  \[
    \texttt{child\_evaluated} = \texttt{calcAttributes ag child\_with\_inh}.
  \]
  This computes the child's own synthesized attributes and checks its local conditions. 

  \item Advance the loop:
  \[
    \texttt{loop (child\_evaluated :: processed) rest occs\_rest}.
  \]
  The measure \(|\texttt{remaining}|\) strictly decreases, ensuring termination. 
\end{enumerate}

When \texttt{remaining} and \texttt{remaining\_o} are both empty, the loop returns \texttt{List.rev processed}, i.e., the fully evaluated children in original 
left-to-right order. Any other shape indicates a mismatch between children and occurrence annotations and triggers an explicit failure. 

\subsection{The Loop Necessity}
The local \texttt{loop} encodes the evaluation discipline that inherited attributes of child \(i\) may depend on the head and on the already evaluated 
siblings \(1..i-1\), but not on later siblings. It does so by constructing \texttt{env\_before} at each step, selecting inherited equations for exactly 
one occurrence, and committing the results before recursing to the next child. This realizes a deterministic, left-to-right strategy and prevents 
accidental dependence on future siblings during inherited-attribute computation. 

\subsection{Relation Back to \texttt{calcAttributes}}
After the \texttt{loop} has produced the fully evaluated children, \texttt{calcAttributes} proceeds to the parent's synthesized pass and condition checking. 
The correctness of those later phases relies on the contract guaranteed by \texttt{loop}: all children are available with their inherited attributes applied 
and their subtrees fully evaluated.


\subsection{\texttt{getRootRule}}
Select the unique grammar rule that matches the current parse-tree node: the rule whose head symbol equals the node's root symbol and whose body 
(sequence of symbols) equals the sequence of the node's immediate children. If the node is a leaf, there is no applicable rule and the function fails. 
If no rule matches, the error message lists candidate rules that share the same head symbol to aid diagnosis. 

Input: the attribute grammar \texttt{ag} and a parse tree \texttt{pt}.  
Output: the matching rule \texttt{rule}.  
Failure cases: leaf nodes, or internal nodes whose children do not match any rule body for the node's head symbol. 

For an internal node, the head symbol is taken from the node's root and the body is the list of child-root symbols. The function searches \texttt{ag.rules} for 
a rule whose head and body match. If none matches, it constructs a readable error with candidate alternatives. 

\begin{verbatim}
let getRootRule (ag: t) (pt: parseTree): rule =
  match pt with
  | Leaf _ ->
      failwith "getRootRule: leaf has no rule"
  | Node (_, children) ->
      let head = getRootSymbol pt in
      let body = List.map getRootSymbol children in
      try Set.find (fun r -> r.head = head && r.body = body) ag.rules with
      | _ ->
          let candidates =
            ag.rules
            |> SetUtil.to_list
            |> List.filter (fun r -> r.head = head)
            |> List.map (fun r ->
                 Printf.sprintf "- %s -> %s"
                   (symb2str r.head)
                   (String.concat " " (List.map symb2str r.body)))
            |> String.concat "\n"
          in
          failwith (Printf.sprintf
            "getRootRule: no rule matches head=%s, body=[%s]\nCandidates with same head:\n%s"
            (symb2str head)
            (String.concat " " (List.map symb2str body))
            (if candidates = "" then "(none)" else candidates))
\end{verbatim}

\subsection{\texttt{apply\_equations\_at\_head}}

Apply a given set of equations to an evaluation environment whose first node must be the head node with a specific symbol. 
Equations are applied in a stable, deterministic order to avoid order-dependent results. The output is a new environment reflecting 
the updates introduced by these equations (e.g., inherited assignments to a specific child, or synthesized assignments to the head). 


Inputs: head symbol \texttt{head\_sym}, a non-empty environment \texttt{nodes} whose first node has symbol \texttt{head\_sym}, and an equation set \texttt{equations}.  
Output: an updated environment \texttt{node list}.  
Failure cases: empty environment or a head-symbol mismatch at the first node. 

The function first verifies that the environment is not empty and that its first node matches \texttt{head\_sym}. It then orders the equations using a 
stable ordering and folds over them, applying each equation with \texttt{eval} to produce the next environment. The fold result is the updated environment. 

\begin{verbatim}
let apply_equations_at_head
    (head_sym : symbol)
    (nodes : node list)
    (equations : equation Set.t)
    : node list =
  let () =
    match nodes with
    | [] ->
        failwith "apply_equations_at_head: empty environment (nodes)"
    | (hs, _) :: _ when hs = head_sym -> ()
    | (hs, _) :: _ ->
        failwith
          (Printf.sprintf
             "apply_equations_at_head: head symbol mismatch (expected %s, got %s)"
             (symb2str head_sym) (symb2str hs))
  in
  let eqs = stable_eq_order head_sym equations in
  List.fold_left
    (fun acc_nodes eq -> eval head_sym eq acc_nodes)
    nodes
    eqs
\end{verbatim}


\subsection{\texttt{check\_conditions\_at\_node}}

Evaluate the rule's boolean conditions at the current node's environment and ensure they all hold. A failing condition aborts evaluation with a message 
identifying the node's head symbol. 

Inputs: head symbol \texttt{head\_sym}, the set of conditions \texttt{conds}, and the current environment \texttt{env}.  
Output: \texttt{unit}.  
Failure cases: any condition evaluates to a non-boolean value or to \texttt{false}. 

Each condition is evaluated with \texttt{evaluate}; the result is coerced with \texttt{as\_bool}. If any evaluation fails to yield \texttt{true}, 
a failure is raised with a diagnostic mentioning the head symbol. 

\begin{verbatim}
let check_conditions_at_node
    (head_sym : symbol)
    (conds : condition Set.t)
    (env : node list)
    : unit =
  Set.iter (fun cond ->
    let v = evaluate head_sym cond env in
    if not (as_bool v) then
      failwith (Printf.sprintf
        "Condition failed at node %s"
        (symb2str head_sym))
  ) conds
\end{verbatim}


\section{The \texttt{accept} function}

The \texttt{accept} receives
a word and an attribute grammar, and returns a boolean indicating whether the word
is accepted by the grammar. The function first converts the attribute grammar into an
equivalent context--free grammar and then relies on the acceptance mechanism of that
grammar.

\begin{verbatim}
let accept (ag: t) (w: word): bool =
    let cfg = ga_to_cfg ag in
    ContextFreeGrammarBasic.accept cfg w
\end{verbatim}



\section{ The \texttt{validate} function}
The \texttt{validate} performs the static checks required before evaluating an attribute grammar on any parse tree. It combines five steps:

\begin{enumerate}
  \item Validate the CFG structure part of the attribute grammar, using the existing CFG module
  \item Type and well-formedness-check all attribute equations,
  \item Ensure all rule conditions are boolean,
  \item Detect dependency cycles among attribute computations. 
\end{enumerate}

If any check fails, an error is raised with an explanatory message. 

The function receives a user-facing name (used in diagnostics) and an attribute-grammar value \texttt{rep}. And returns normally if and only 
if all checks pass. Otherwise it reports an error with \texttt{Error.error} or raises a failure from helper routines. 

The grammar is first projected to a CFG with \texttt{ga\_to\_cfg}. This preserves terminals, variables, initial symbol, and the bare head\(\to\)body production 
structure. The resulting CFG is then validated by \texttt{ContextFreeGrammarPrivate.validate}, which checks basic well-formedness (e.g., symbols belong to the 
declared sets, rules reference known symbols, etc.). 

\texttt{validateEquations} traverses all rules and their equations. For each equation \( \texttt{lhs} = \texttt{rhs} \):
\begin{itemize}
  \item the left-hand side must be an \texttt{Apply(attr,(var,i))} to a valid target (head or a concrete child occurrence in range); the variable must be 
  a declared grammar variable; and the direction is checked: assignments to the head require synthesized attributes, and assignments to children require 
  inherited attributes; 
  \item both sides are type-checked via \texttt{validateExp}, ensuring equal, non-\texttt{"error"} types. Operators are checked for arity and operand types 
  (e.g., integer arithmetic, string concatenation, comparisons). 
\end{itemize}
Any mismatch (invalid target, out-of-range occurrence, wrong direction, unknown symbol/attribute, or type error) triggers an error at this stage. 

\texttt{validateConditions} ensures every condition expression attached to a rule statically type-checks as boolean (again using \texttt{validateExp}). 
Non-boolean conditions are rejected with a diagnostic. 

\texttt{has\_cycles} constructs a directed dependency graph whose nodes are attribute occurrences identified by 
\((\textit{symbol}, \textit{occurrence}, \textit{attribute})\). For each equation, edges are added from every attribute reference in the right-hand side to 
the left-hand side target, within the context of the same rule. A depth-first search detects any cycle; if present, returns false, and add the error message 
\enquote{Attribute dependency cycle detected} to the errors list. This prevents unschedulable attribute computations at evaluation time. 

\begin{verbatim}
let validate (name: string) (rep: t): unit =
  let cfg = ga_to_cfg rep in
  ContextFreeGrammarPrivate.validate name cfg;
  validateEquations name rep;
  validateConditions name rep;
  if has_cycles rep then
    Error.error name "Attribute dependency cycle detected" ()
\end{verbatim}


\subsection{Auxiliary Methods Used by \texttt{validate}}

The following helpers iterate over all rules and check each equation or condition individually:

\begin{verbatim}
let validateEquations (name: string) (rep: t): unit =
  Set.iter (fun r ->
    Set.iter (fun eq ->
      validateEquation name rep eq r
    ) r.equations
  ) rep.rules

let validateConditions (name: string) (rep: t): unit =
  Set.iter (fun r ->
    Set.iter (fun eq ->
      validateCondition name rep eq r
    ) r.conditions
  ) rep.rules
\end{verbatim}
These use rule-scoped checkers: \texttt{validateEquation} confirms target well-formedness, direction (inherited vs.\ synthesized), variable membership, 
and type equality between left and right expressions; \texttt{validateCondition} verifies that the condition's static type is boolean. 


The cycle check reduces to a graph problem:
\begin{verbatim}
let has_cycles (ag : t) : bool =
  ag |> build_dep_graph |> has_cycle_graph
\end{verbatim}
Here \texttt{build\_dep\_graph} indexes each attribute occurrence into a node id and collects edges from RHS references to LHS targets for every equation 
of every rule. \texttt{has\_cycle\_graph} runs a DFS-based cycle test over the adjacency lists. If any back-edge is found, the function returns \texttt{true}. 


\section{Cycle Detection}
Attribute dependencies are modelled as a directed graph: each node represents a concrete attribute occurrence 
\((\textit{symbol}, \textit{occurrence}, \textit{attribute})\), and each equation adds edges from every attribute referenced on the right-hand 
side to the left-hand side target. A DFS-based test on this graph reports whether any cycle exists. If a cycle is found, validation fails. 

\subsection{\texttt{build\_dep\_graph}}
this function constructs the dependency graph of the whole grammar and return it as adjacency lists. 

\begin{itemize}
  \item Map each occurrence \((s, k, a)\) to a unique node id (hash table). 
  \item For each rule equation \(LHS = RHS\): translate \(LHS\) to its destination node; collect all attribute refs in \(RHS\) and add edges \(\text{ref} \to LHS\). 
  \item Materialize the edge set into an \texttt{int list array}. 
\end{itemize}

\begin{verbatim}
type node_key = symbol * int * attribute

let build_dep_graph (ag : t) : int list array =
  let index_tbl : (node_key, int) Hashtbl.t = Hashtbl.create 64 in
  let counter = ref 0 in
  let get_id key =
    match Hashtbl.find_opt index_tbl key with
    | Some id -> id
    | None ->
        let id = !counter in
        incr counter;
        Hashtbl.add index_tbl key id;
        id
  in
  let edges = ref [] in
  let rules = SetUtil.to_list ag.rules in
  List.iter (fun r ->
    Set.iter (fun eq ->
      match eq with
      | Apply (attrL, (vL,iL)), rhs ->
          let (sL,jL) = sym_occ_of_ref r (vL,iL) in
          let dst = get_id (sL, jL, attrL) in
          let refs = refs_in_expr r rhs in
          List.iter (fun (s,j,a) ->
            let src = get_id (s,j,a) in
            edges := (src, dst) :: !edges
          ) refs
      | _ -> ()
    ) r.equations
  ) rules;
  let n = !counter in
  let adj = Array.make n [] in
  List.iter (fun (u,v) -> adj.(u) <- v :: adj.(u)) !edges;
  adj
\end{verbatim}

\subsection{\texttt{has\_cycle\_graph}}
This function decides if the directed graph contains any cycle using DFS with a recursion stack. 

Maintain two boolean arrays: \texttt{visited} and \texttt{stack}. A back-edge to a node currently in \texttt{stack} signals a cycle; otherwise, explore all 
successors. 

\begin{verbatim}
let rec dfs adj u visited stack =
  if stack.(u) then true
  else if visited.(u) then false
  else (
    visited.(u) <- true;
    stack.(u) <- true;
    let found = List.exists (fun v -> dfs adj v visited stack) adj.(u) in
    stack.(u) <- false;
    found
  )

let has_cycle_graph adj =
  let n = Array.length adj in
  let visited = Array.make n false in
  let stack = Array.make n false in
  let rec loop i =
    if i = n then false
    else if (not visited.(i)) && dfs adj i visited stack
    then true
    else loop (i+1)
  in
  loop 0
\end{verbatim}
